% ======================================================================
% preamble.tex -- packages, theorem environments, draft-mode machinery
% ======================================================================

% ---------------------------------------------------------------- fonts
% (guarded so the skeleton also compiles under XeTeX/LuaTeX engines,
% e.g. tectonic; pdflatex behavior is unchanged)
\ifdefined\XeTeXrevision\else\ifdefined\directlua\else
  \usepackage[utf8]{inputenc}
  \usepackage[T1]{fontenc}
\fi\fi
\IfFileExists{microtype.sty}{\usepackage{microtype}}{}

% ------------------------------------------------------------- packages
\usepackage{booktabs}
\usepackage{graphicx}
\usepackage{amsmath}
\usepackage{amssymb}
\usepackage{amsthm}
\usepackage{xcolor}
\usepackage{colortbl}
\usepackage{xspace}
\usepackage{natbib}          % official ICLR style ships its own .bst; see main.tex

% Boxed-theory environment: tcolorbox if available, framed as fallback.
\IfFileExists{tcolorbox.sty}{%
  \usepackage{tcolorbox}%
  \newtcolorbox{theorybox}{%
    colback=gray!6, colframe=black!70, boxrule=0.6pt, arc=1.5mm,
    left=7pt, right=7pt, top=6pt, bottom=6pt}%
}{%
  \usepackage{framed}%
  \newenvironment{theorybox}{\begin{framed}}{\end{framed}}%
}

% hyperref + cleveref last
\usepackage[colorlinks=true,
            linkcolor=blue!55!black,
            citecolor=green!40!black,
            urlcolor=blue!55!black]{hyperref}
\usepackage[capitalize,noabbrev]{cleveref}

% --------------------------------------------------- theorem environments
\theoremstyle{plain}
\newtheorem{theorem}{Theorem}
\newtheorem{proposition}{Proposition}
\newtheorem{lemma}{Lemma}
\newtheorem{corollary}{Corollary}
\theoremstyle{definition}
\newtheorem{definition}{Definition}
\newtheorem{assumption}{Assumption}
\theoremstyle{remark}
\newtheorem{remark}{Remark}

% Fixed-label statements for the boxed theory blocks (T-numbered, so they
% survive reordering of the ordinary theorem counter).
\newtheorem*{propositionT}{Proposition~T1}
\newtheorem*{corollaryT}{Corollary~T2}
\newtheorem*{propositionTrank}{Proposition~T3}

% ======================================================================
% Draft-mode machinery
% ======================================================================

% ----- frozen-results toggle ------------------------------------------
% While results are UNFROZEN (handoff-sourced numbers, not yet regenerated
% from data/frozen/), every number referenced via \unfrozen{...} carries a
% gray "UNFROZEN" superscript. Flip the boolean below once
% scripts/make_tables.py has regenerated results_macros.tex from
% data/frozen/ -- all markers disappear and numbers render plainly.
\newif\ifFrozenResults
\FrozenResultstrue   % results frozen 2026-08-18 (manifest in data/frozen/MANIFEST.json)

\newcommand{\unfrozen}[1]{%
  \ifFrozenResults
    {#1}%
  \else
    {#1}\textsuperscript{\,\textcolor{gray}{\tiny\textsf{UNFROZEN}}}%
  \fi}

% ----- researcher TODO margin notes -----------------------------------
% \todoresearcher{...}       margin note (body text only)
% \todoresearcherinline{...} inline note -- REQUIRED inside floats and
%                            captions, where \marginpar is illegal.
\newcommand{\todoresearcher}[1]{%
  \marginpar{\raggedright\scriptsize\textcolor{red!70!black}{%
    \textbf{TODO(researcher):} #1}}}
\newcommand{\todoresearcherinline}[1]{%
  \textcolor{red!70!black}{\scriptsize\textbf{TODO(researcher):} #1}}

% ----- misc helper macros ---------------------------------------------
\newcommand{\method}{TopoSHAP\xspace}
\DeclareMathOperator*{\argmax}{arg\,max}

% Bolded one-sentence claim opening an experiment subsection.
\newcommand{\expclaim}[1]{\noindent\textbf{#1}\par\smallskip}

% Boxed practitioner's recipe: plain-language numbered steps, no math.
\newenvironment{recipebox}{\begin{theorybox}\noindent
  \textbf{Practitioner's recipe.}\begin{enumerate}\setlength{\itemsep}{1pt}}
  {\end{enumerate}\end{theorybox}}

% Neighborhood names (message-passing routes); always "neighborhoods",
% never "relations".
\newcommand{\nbhd}[1]{\texttt{#1}}

% Placeholder box for figures whose PDF has not been generated yet.
\newcommand{\figureplaceholder}[1]{%
  \fbox{\parbox[c][0.30\linewidth][c]{0.95\linewidth}{%
    \centering\color{gray}\textsf{FIGURE PLACEHOLDER}\\[2pt]
    \texttt{\detokenize{#1}}}}}

% Convenience shorthand
\newcommand{\eg}{e.g.,\xspace}
\newcommand{\ie}{i.e.,\xspace}
\newcommand{\rk}{\operatorname{rk}}
